\chapter{Justificación}
\label{cap:justificacion}

\section{Justificación científica}

El seguimiento de un precio de referencia por un AMM ya se ha estudiado desde el análisis formal y la simulación \cite{angeris2020uniswap}, y también la rentabilidad de la provisión de liquidez frente a un CEX \cite{campbell2025optimal,milionis2022lvr}. Lo que este trabajo añade no es el fenómeno sino la forma de examinarlo: separa el efecto de la capacidad de arbitraje y el de la reacción de los proveedores mediante comparaciones pareadas con flujos de órdenes comunes, y contrasta una descripción reducida de dos variables con el simulador del que proviene.

La relevancia científica está en la delimitación de lo que cada representación puede afirmar. Una identidad analítica, como el aumento del impacto al disminuir las reservas, no requiere experimentación; una afirmación sobre arbitraje limitado o retiro adaptativo sí, porque depende de reglas de comportamiento y de un horizonte. Distinguir ambos niveles evita presentar como hallazgo empírico una consecuencia algebraica y evita también descartar como trivial un efecto que solo aparece bajo ciertas profundidades. El resultado del capítulo~\ref{cap:resultados} ilustra el punto: el efecto de la liquidez adaptativa sobre la respuesta máxima no es uniforme, sino que depende de la profundidad ensayada.

Además, la formulación como contraste confirmatorio, con hipótesis, métricas y criterios de rechazo fijados antes de simular las semillas de confirmación, responde a una crítica frecuente a los ABM: que se ajustan a posteriori hasta mostrar el fenómeno deseado \cite{fagiolo2007critical}. El protocolo no elimina esa crítica, pero hace explícito qué se decidió antes de observar los resultados y qué no. La señal de que el protocolo es útil es que también documenta resultados desfavorables a la expectativa inicial: el mapa reducido no supera a un autorregresivo escalar, y la precisión objetivo para H1 no se alcanzó con el número de semillas previsto.

\section{Justificación metodológica}

Una comparación sencilla entre un escenario con arbitraje y otro sin él no permite atribuir la diferencia al arbitraje si los flujos de órdenes difieren entre ambos. La comparación pareada, con los mismos choques exógenos y las mismas órdenes de ruido pregeneradas, aísla la regla modificada, y es la lógica que Palit \cite{palit2022shuffled} aplica al reproducir eventos históricos bajo un agente sin inteligencia. La inferencia con la semilla como unidad independiente evita tratar como muestras distintas las celdas de un mismo flujo, y el control de Holm limita el riesgo de falsos positivos al evaluar varias hipótesis a la vez.

El uso de configuraciones y semillas reservadas para evaluar el mapa reducido responde a un problema conocido en la identificación de modelos: el ajuste dentro de la muestra sobrestima la capacidad de descripción \cite{lamperti2018surrogates}. Ajustar el mapa con un conjunto de exploración y evaluarlo en otro permite distinguir entre una descripción que generaliza y una que solo reproduce las trayectorias de las que se estimó. La comparación contra la persistencia y contra un modelo autorregresivo escalar establece qué aporta el segundo estado del mapa. Sin esas referencias, un error pequeño no sería interpretable, porque la persistencia ya es una predicción competitiva a un paso.

Finalmente, el registro de cada ejecución con identificador, semilla, parámetros y verificación de invariantes \cite{grimm2020odd,railsback2019agent} permite que otra persona repita el experimento y compruebe las cifras publicadas. La reproducibilidad completa del experimento confirmatorio (identificador de ejecución, sumas de verificación y comando) se documenta en el anexo~\ref{cap:anexos}.

\section{Justificación aplicada y pedagógica}

Los resultados no están destinados a prescribir estrategias de inversión ni a calibrar un protocolo real. Su utilidad es de otro orden. Primero, para el diseño de mercados descentralizados, muestran que la profundidad inicial y la capacidad de arbitraje condicionan la rapidez de corrección, y que la reacción de los proveedores solo amplifica una perturbación bajo ciertas condiciones; esa información orienta qué parámetros conviene observar en un análisis posterior con datos. Segundo, en la formación de ingenieros de sistemas, el micromundo ofrece un caso en que la relación entre mecanismos locales, restricciones y comportamiento agregado puede explorarse con un código pequeño y verificable. Tercero, el sitio web asociado al trabajo reproduce el escenario con una interfaz interactiva, pensada para estudiar de manera intuitiva los efectos de reducir la profundidad o de limitar el arbitraje; su valor pedagógico no se ha evaluado con estudiantes, por lo que se menciona como línea de trabajo futuro (capítulo~\ref{cap:futuro}).

Los límites de estas aplicaciones están declarados: el modelo es sintético, no emplea datos de mercado y representa un único par con un AMM de producto constante. Cualquier uso fuera de ese dominio exige evidencia adicional.
